% Part: proof-theory
% Chapter: sequent-calculus
% Section: proof-examples

\documentclass[../../../include/open-logic-section]{subfiles}

\begin{document}

\olfileid{pt}{seq}{ex}

\olsection{સાબિતીઓનાં ઉદાહરણો}

સિક્વન્ટોની !!{proof}s આપવા
સામાન્ય રીતે અંતિમ સિક્વન્ટથી ઉપર તરફ કામ કરવું અનુકૂળ છે:
તેમાંના !!a{formula}ને મુખ્ય !!{formula} તરીકે પસંદ કરી,
અનુરૂપ નિયમના આધારો સિક્વન્ટની ઉપર લખીએ;
પછી સ્વયંસિદ્ધો સુધી પહોંચીએ ત્યાં સુધી પ્રક્રિયા દોહરાવીએ.
\sourcecorrection{OLMD-315}{મૂળ પાછી તરફ સાબિતી-શોધમાં નિષ્કર્ષમાંથી પસંદ કરેલા સૂત્રને સક્રિય કહે છે. નિયમની વ્યાખ્યા મુજબ નિષ્કર્ષનું પસંદ કરેલું સૂત્ર મુખ્ય છે; આધારોમાં તેનાથી મળતાં સૂત્રો સક્રિય છે. અહીં શરૂઆત અને બીજા અનુસરણના પગલામાં મુખ્ય સૂત્ર કહ્યું છે.}
દાખલા તરીકે આ સિક્વન્ટ સાબિત કરવો છે એમ ધારો:
\[(!C \land !D) \lif !E \Sequent !C \lif (!D \lif !E)\]
$!C \lif (!D \lif !E)$ વિચારીએ:
તે $!A \lif !B$ સ્વરૂપનું છે
અને સિક્વન્ટની જમણી બાજુ છે.
આખા સિક્વન્ટને \Log{G1c}ના $\RightR{\lif}$ નિયમ સાથે સરખાવીએ,
\[\Axiom$ !A, \Gamma \fCenter \Delta, !B$
\RightLabel{\RightR{\lif}}
\UnaryInf$ \Gamma \fCenter \Delta, !A \lif !B $
\DisplayProof\]
તો $\Gamma = \{(!C \land !D) \lif !E\}$,
$\Delta = \emptyset$, $!A \ident !C$
અને $!B \ident !D \lif !E$ લેવા જોઈએ તે સૂચવે છે.
\sourcecorrection{OLMD-316}{મૂળ આ પગલામાં સંદર્ભને નિષ્કર્ષના જમણા સૂત્રથી ભરે છે. દર્શાવેલા સિક્વન્ટનો પૂર્વાંગ તો સંયોજનમાંથી અનુસરણ છે; Gammaમાં એ જ પૂર્વાંગ રાખ્યો છે, જેમ તરત પછીનું વૃક્ષ કરે છે.}
એટલે સાબિતીનું છેલ્લું અનુમાન આ હોઈ શકે:
\[\Axiom$ !C, (!C \land !D) \lif !E \fCenter !D \lif !E$
\RightLabel{\RightR{\lif}}
\UnaryInf$(!C \land !D) \lif !E \fCenter !C \lif (!D \lif !E)$
\DisplayProof\]
$!D \lif !E$ને મુખ્ય સૂત્ર ગણીને વિચાર દોહરાવીએ,
તો આ મળે:
\[\Axiom$ !D, !C, (!C \land !D) \lif !E \fCenter !E$
\RightLabel{\RightR{\lif}}
\UnaryInf$!C, (!C \land !D) \lif !E \fCenter !D \lif !E$
\RightLabel{\RightR{\lif}}
\UnaryInf$(!C \land !D) \lif !E \fCenter !C \lif (!D \lif !E)$
\DisplayProof\]
હવે $(!C \land !D) \lif !E$ પર ધ્યાન આપવું
અને \Log{G1c}નો $\LeftR{\lif}$ નિયમ વાપરવો જોઈએ:
\[\Axiom$ \Gamma \fCenter \Delta, !A$
\Axiom$ !B, \Gamma \fCenter \Delta$
\RightLabel{\LeftR{\lif}}
\BinaryInf$ !A \lif !B, \Gamma \fCenter \Delta$
\DisplayProof\]
તેથી આ મળે:
\[\Axiom$!D, !C \fCenter !E, !C \land !D$
\Axiom$!D, !C, !E \fCenter !E$
\RightLabel{\LeftR{\lif}}
\BinaryInf$ !D, !C, (!C \land !D) \lif !E \fCenter !E$
\RightLabel{\RightR{\lif}}
\UnaryInf$!C, (!C \land !D) \lif !E \fCenter !D \lif !E$
\RightLabel{\RightR{\lif}}
\UnaryInf$(!C \land !D) \lif !E \fCenter !C \lif (!D \lif !E)$
\DisplayProof\]
સૌથી ઉપરના ડાબા સિક્વન્ટમાં !!{formula} $!C \land !D$ને
મુખ્ય ગણી $\RightR{\land}$ લગાવીએ, તો આ મળે:
\[
\Axiom$!D, !C \fCenter !E, !C$
\Axiom$!D, !C \fCenter !E, !D$
\RightLabel{\RightR{\land}}
\BinaryInf$!D, !C \fCenter !E, !C \land !D$
\Axiom$!D, !C, !E \fCenter !E$
\RightLabel{\LeftR{\lif}}
\BinaryInf$ !D, !C, (!C \land !D) \lif !E \fCenter !E$
\RightLabel{\RightR{\lif}}
\UnaryInf$!C, (!C \land !D) \lif !E \fCenter !D \lif !E$
\RightLabel{\RightR{\lif}}
\UnaryInf$(!C \land !D) \lif !E \fCenter !C \lif (!D \lif !E)$
\DisplayProof
\]
\sourcecorrection{OLMD-317}{મૂળના આ વૃક્ષમાં બે આધારવાળા અનુસરણ-ડાબા પગલાનું લેબલ અનુસરણ-જમણું છે. બંને આધાર અને મુખ્ય પૂર્વાંગ પ્રમાણે તેને ડાબું કર્યું છે.}
અત્યાર સુધી બધું યોગ્ય છે.
અત્યાર સુધી વાપરેલા નિયમો \Log{G1c} અને~\Log{G3c}માં સમાન હોવાથી
આ બધું \Log{G3c}માં પણ ચાલે છે.
આપણે !!a{proof}ના એવા અંશ સુધી પહોંચ્યા છીએ
જેના દરેક સૌથી ઉપરના સિક્વન્ટમાં ડાબે અને જમણે એ જ !!{formula} છે.
પરંતુ તેઓ હજી સ્વયંસિદ્ધો નથી.

\Log{G1c}માં $!A \Sequent !A$ સ્વરૂપના સ્વયંસિદ્ધોમાંથી
સૌથી ઉપરના સિક્વન્ટોની !!{proof}s આપવી પડે.
શિથિલીકરણના નિયમોથી આ સરળતાથી થઈ શકે.
દાખલા તરીકે ડાબે સૌથી ઉપરનો સિક્વન્ટ
$!D, !C \Sequent !E, !C$ આ રીતે !!{prove}d થઈ શકે:
\[
\Axiom$!C \fCenter !C$
\RightLabel{\LeftR{\Weakening}}
\UnaryInf$!D, !C \fCenter !C$
\RightLabel{\RightR{\Weakening}}
\UnaryInf$!D, !C \fCenter !E, !C$
\DisplayProof
\]
\sourcecorrection{OLMD-318}{મૂળ આ નાનાં શિથિલીકરણ-વૃક્ષ પછી તેને દર્શાવતો DisplayProof આદેશ છોડે છે. માત્ર તે આદેશ ઉમેર્યો છે; વૃક્ષનાં સૂત્રો અને નિયમો અચળ છે.}
આખરે \Log{G1c}માં આ સ્વરૂપની !!a{proof} મળે:
\[
\Axiom$!C \fCenter !C$
\RightLabel{\LeftR{\Weakening}}
\UnaryInf$!D, !C \fCenter !C$
\RightLabel{\RightR{\Weakening}}
\UnaryInf$!D, !C \fCenter !E, !C$
\Axiom$!D \fCenter !D$
\RightLabel{\LeftR{\Weakening}}
\UnaryInf$!D, !C \fCenter !D$
\RightLabel{\RightR{\Weakening}}
\UnaryInf$!D, !C \fCenter !E, !D$
\RightLabel{\RightR{\land}}
\BinaryInf$!D, !C \fCenter !E, !C \land !D$
\Axiom$!E \fCenter !E$
\RightLabel{\LeftR{\Weakening}}
\UnaryInf$!C, !E \fCenter !E$
\RightLabel{\LeftR{\Weakening}}
\UnaryInf$!D, !C, !E \fCenter !E$
\RightLabel{\LeftR{\lif}}
\BinaryInf$ !D, !C, (!C \land !D) \lif !E \fCenter !E$
\RightLabel{\RightR{\lif}}
\UnaryInf$!C, (!C \land !D) \lif !E \fCenter !D \lif !E$
\RightLabel{\RightR{\lif}}
\UnaryInf$(!C \land !D) \lif !E \fCenter !C \lif (!D \lif !E)$
\DisplayProof
\]
\sourcecorrection{OLMD-319}{મૂળના જમણા સ્વયંસિદ્ધમાંથી પહેલા શિથિલીકરણના પગલામાં ઉત્તરાંગ બદલીને બીજું સૂત્ર કરે છે. અહીં ઉત્તરાંગનું એ જ $!E$ રાખ્યું છે; શિથિલીકરણ ફક્ત ડાબે $!C$ ઉમેરે છે.}
આ સાબિતીની સંરચના (ખાસ કરીને ઊંચાઈ)
!!{formula}s $!C$, $!D$ અને~$!E$ કયા છે તેનાથી સ્વતંત્ર છે.
ઉપરની !!{proof}માં આ ચિહ્નોની જગ્યાએ
કોઈ પણ !!{formula}s મૂકી શકાય,
અને પરિણામ હજી~\Log{G1c}ની !!{proof} હોય.
\Log{G1c}ની સાબિતી \emph{ઢાંચારૂપ} છે.

\Log{G3c}માં સ્વયંસિદ્ધો
$!A, \Gamma \Sequent \Delta, !A$ સ્વરૂપના સિક્વન્ટો છે,
જ્યાં $!A$ આણ્વિક હોવું જોઈએ.
એટલે $!D, !C \Sequent !E, !C$ એ \Log{G3c}ના
સ્વયંસિદ્ધ જેવો લાગે છે
($\Gamma = \{!D\}$; $\Delta = \{!E\}$),
પરંતુ $!C$ ખરેખર આણ્વિક હોય ત્યારે
તે આ મેળ ખાતી આવૃત્તિઓથી \emph{ખરેખર} સ્વયંસિદ્ધ છે.
$!C$, $!D$ અને $!E$ આણ્વિક !!{formula}s હોય,
તો આપણો સાબિતી-અંશ \Log{G3c}માં સંપૂર્ણ સાબિતી છે.
\sourcecorrection{OLMD-320}{મૂળ આણ્વિકતાને આ સિક્વન્ટો સ્વયંસિદ્ધ હોવાની આવશ્યક શરત કહે છે. અહીં બતાવેલી મેળ ખાતી આવૃત્તિઓ આણ્વિક હોય તે પૂરતી શરત છે; બીજા આણ્વિક સૂત્રનો બંને બાજુ મેળ અથવા ડાબે અસત્ય હોવાથી પણ સિક્વન્ટ સ્વયંસિદ્ધ હોઈ શકે. તેથી પૂરતી શરત તરીકે કહ્યું છે.}

સખત રીતે જોતાં આપણે
\[\Log{G3c} \Proves (!C \land !D) \lif !E \fCenter !C \lif (!D \lif
!E),\]
બતાવ્યું નથી, છતાં વ્યવહારમાં આપણે આટલું જ કરવાની
જરૂર છે (અને આટલું જ કરી શકીએ).
કારણ કે $!A$ પોતે આણ્વિક ન હોય તો પણ
$!A, \Gamma \Sequent \Delta, !A$ સ્વરૂપનો દરેક સિક્વન્ટ
\Log{G3c}માં સાબિત થઈ શકે છે.
$!A$ની ઊંડાઈ~$\depth{!A}$ પર આગમનથી આ બતાવી શકીએ.
\sourcecorrection{OLMD-321}{મૂળ આ વાક્યના પ્રદર્શનને G1c કહે છે, પરંતુ એ તંત્રની સંપૂર્ણ ઢાંચારૂપ સાબિતી ઉપર આપી ચૂક્યું છે. અહીં હજી પૂર્ણ ન થયેલી વાત G3cમાં અણઆણ્વિક સ્વયંસિદ્ધોના સ્થાને સાબિતીઓ આપવાની છે; તેથી તંત્રનું નામ G3c કર્યું છે.}

\begin{defn}\ollabel{defn:depth}
!!a{formula}ની \emph{ઊંડાઈ}~$\depth{!A}$ આ રીતે વ્યાખ્યાયિત છે:
\begin{enumerate}
  \item $!A$ આણ્વિક હોય, તો $\depth{!A} = 0$.
  \item $!A \ident \lnot !B$, $!A \ident \lforall[x][!B]$
  અથવા $!A \ident \lexists[x][!B]$ હોય,
  તો $\depth{!A} = \depth{!B} + 1$.
  \item $!A \ident (!B \land !C)$, $(!B \lor !C)$
  અથવા $(!B \lif !C)$ હોય,
  તો $\depth{!A} = \max(\depth{!B},\depth{!C}) + 1$.
\end{enumerate}
\sourcecorrection{OLMD-322}{મૂળ પરિમાણકનાં બંને કિસ્સામાં અંદરનું સૂત્ર $A$ લખે છે, પરંતુ ઊંડાઈની જમણી બાજુ અસંબંધિત $!B$ છે. બંને અંદરના સૂત્રોને એ જ $!B$ કર્યા છે, જેમ નિષેધના કિસ્સામાં છે.}
\end{defn}

કોઈ પણ પદ~$t$ માટે $\depth{A(x)} = \depth{A(t)}$
સાબિત કરવું મુશ્કેલ નથી.
આપણા માટે મહત્ત્વની વાત એ છે કે વિઘટનના તાર્કિક નિયમોમાં
મુખ્ય !!{formula}ની ઊંડાઈ તેમાંથી મળતાં સક્રિય
ઉપ-!!{formula}sની ઊંડાઈ કરતાં મોટી છે.
\sourcecorrection{OLMD-323}{મૂળ આ ઊંડાઈની કડક અસમાનતા દરેક નિયમ માટે કહે છે. સંકોચન જેવા બંધારણીય નિયમમાં એક જ સૂત્રની ઊંડાઈ બદલાતી નથી; મુખ્ય સૂત્રની નકલ જાળવતા પરિમાણક-નિયમમાં પણ એ નકલની ઊંડાઈ સમાન છે. અહીં ઊંડાઈ ઘટતી હોય તે સાચા વિઘટિત ઉપસૂત્રોની વાત સ્પષ્ટ કરી છે.}
દાખલા તરીકે:
\begin{align*}
\depth{P(x)} = \depth{P(a)} & = 0,\\
\depth{\lforall[x][P(x)]} & = 1,\\
\depth{(\lforall[x][P(x)] \lif P(a))} & = \max(1,0) + 1 = 2.
\end{align*}
આથી~$\depth{!A}$ પર આગમનથી પરિણામો સાબિત કરી શકીએ,
જેમ કે આ પરિણામ.

\begin{prop}\ollabel{prop:G1c-axioms}
કોઈ પણ~$!A$ માટે $!A \Sequent !A$ની $ \Log{G1c}$માં
એવી !!a{proof} છે જેના બધા સ્વયંસિદ્ધો આણ્વિક છે.
\end{prop}

\begin{proof}
  $\depth{!A}$ પર આગમનથી સાબિતી કરીએ.
  $\depth{!A} = 0$ હોય, તો $!A$ આણ્વિક છે
  અને $!A \Sequent !A$ પોતે~\Log{G1c}ની
  જરૂરી સ્વરૂપની !!{proof} છે.
  $\depth{!A} > 0$ હોય, તો $!A$ ઓછી ઊંડાઈનાં
  !!{formula}sમાંથી બનેલું છે,
  દાખલા તરીકે $!A \ident \lnot !B$,
  $!A \ident (!B \land !C)$
  અથવા $!A \ident \lforall[x][A(x)]$.
  આગમન-ધારણા એ છે કે
  $\depth{!D} < \depth{!A}$ હોય તે બધાં $!D$ માટે
  આણ્વિક સ્વયંસિદ્ધોમાંથી $\Log{G1c} \Proves !D \Sequent !D$ છે.

  $!A \ident \lnot !B$ હોય, તો $\depth{!B} < \depth{!A}$ છે.
  આગમન-ધારણાથી \Log{G1c}માં આણ્વિક સ્વયંસિદ્ધોમાંથી
  $!B \Sequent !B$ની !!a{proof}~$\pi_1$ છે.
  તેને આ રીતે $ \lnot !B \Sequent \lnot !B$ની !!a{proof}માં વિસ્તારી શકીએ:
\[
  \AxiomC{}
  \RightLabel{$\pi_1$}
  \Deduce$!B \fCenter !B$
  \RightLabel{\LeftR{\lnot}}
  \UnaryInf$\lnot !B, !B \fCenter $
  \RightLabel{\RightR{\lnot}}
  \UnaryInf$\lnot !B \fCenter \lnot !B$
  \DisplayProof
\]

  $!A \ident (!B \land !C)$ હોય,
  તો $\depth{!B} < \depth{!A}$ અને $\depth{!C} < \depth{!A}$ છે.
  આગમન-ધારણાથી \Log{G1c}માં આણ્વિક સ્વયંસિદ્ધોમાંથી
  $!B \Sequent !B$ અને $!C \Sequent !C$ની
  !!{proof}s~$\pi_1$ અને~$\pi_2$ છે.
  તેમને $!B \land !C \Sequent !B \land !C$ની
  !!a{proof}માં આ રીતે વિસ્તારી શકીએ:
\[
  \AxiomC{}
  \RightLabel{$\pi_1$}
  \Deduce$!B, \fCenter !B$
  \RightLabel{\LeftR{\Weakening}}
  \UnaryInf$!B, !C, \fCenter !B$
  \RightLabel{\LeftR{\land}}
  \UnaryInf$!B \land !C, !C\fCenter !B$
  \RightLabel{\LeftR{\land}}
  \UnaryInf$!B \land !C, !B \land !C \fCenter !B$
  \AxiomC{}
  \RightLabel{$\pi_2$}
  \Deduce$!C, \fCenter !C$
  \RightLabel{\LeftR{\Weakening}}
  \UnaryInf$!B, !C, \fCenter !C$
  \RightLabel{\LeftR{\land}}
  \UnaryInf$!B \land !C, !C \fCenter !C$
  \RightLabel{\LeftR{\land}}
  \UnaryInf$!B \land !C, !B \land !C \fCenter !C$
  \RightLabel{\RightR{\land}}
  \BinaryInf$!B \land !C, !B \land !C \fCenter !B \land !C$
  \RightLabel{\LeftR{\Contraction}}
  \UnaryInf$!B \land !C \fCenter !B \land !C$
  \DisplayProof
  \]
  \sourcecorrection{OLMD-324}{મૂળના છેલ્લાં બે પગલામાં સંયોજન-જમણાનો નિષ્કર્ષ એક ડાબી આવૃત્તિ ખોટી રીતે ગુમાવે છે અને પછી ડાબું સંકોચન જમણું સંયોજન બદલી નાખે છે. પહેલા પગલામાં ડાબે બે સંયોજન રાખ્યાં છે; સંકોચન પછી ડાબે એક સંયોજન અને જમણે એ જ સંયોજન રહે છે.}

  હવે $!A \ident \lforall[x][!B(x)]$ છે એમ ધારો.
  $\lforall[x][!B(x)]$માં ન આવતો !!a{constant}~$c$ લો.
  ત્યારે $\depth{!B(c)} = \depth{!B(x)} < \depth{!A}$ છે.
  આગમન-ધારણાથી \Log{G1c}માં આણ્વિક સ્વયંસિદ્ધોમાંથી
  $!B(c) \Sequent !B(c)$ની !!a{proof}~$\pi_1$ છે.
  તેને $\lforall[x][!B(x)] \Sequent \lforall[x][!B(x)]$ની
  !!a{proof}માં આ રીતે વિસ્તારી શકીએ:
\[
  \AxiomC{}
  \RightLabel{$\pi_1$}
  \Deduce$!B(c) \fCenter !B(c)$
  \RightLabel{\LeftR{\lforall}}
  \UnaryInf$\lforall[x][!B(x)] \fCenter !B(c)$
  \RightLabel{\RightR{\lforall}}
  \UnaryInf$\lforall[x][!B(x)] \fCenter \lforall[x][!B(x)]$
  \DisplayProof
\]
  \sourcecorrection{OLMD-325}{મૂળ આ પરિચ્છેદના કેટલાક સૂત્રોમાં સામાન્ય $B$ અને કેટલાકમાં સૂત્ર-સ્થાનધારક $!B$ વાપરે છે; વૃક્ષમાં $!B$ છે. ગદ્યની સર્વપરિમાણક-ઓળખ અને દાખલની ઊંડાઈ પણ એ જ $!B$ સાથે સુસંગત કરી છે.}

  બાકીના કિસ્સા સ્વાધ્યાય તરીકે છોડ્યા છે.
\end{proof}

\begin{prob}
  $!A \ident (!B \lor !C)$, $!A \ident (!B \lif !C)$
  અને $!A \ident \lexists[x][B(x)]$ હોય તે કિસ્સા કરીને
  \olref{prop:G1c-axioms}ની સાબિતી પૂર્ણ કરો.
\end{prob}

આપણા હેતુ માટે $!A \ident \lnot !B$ના કિસ્સામાં
આ !!{proof} પણ વાપરી શક્યા હોત:
\[
\AxiomC{}
\RightLabel{$\pi_1$}
\Deduce$!B \fCenter !B$
\RightLabel{\RightR{\lnot}}
\UnaryInf$\fCenter !B, \lnot !B$
\RightLabel{\LeftR{\lnot}}
\UnaryInf$\lnot !B \fCenter \lnot !B$
\DisplayProof
\]
પરંતુ કેટલાંક સિક્વન્ટ તંત્રોમાં આ ન ચાલ્યું હોત.
અંતઃપ્રજ્ઞાવાદી સિક્વન્ટ કલન \Log{G1i} એ \Log{G1c}ને અનુરૂપ છે;
તેમાં દરેક સિક્વન્ટની જમણી બાજુ વધુમાં વધુ એક !!{formula}ની મર્યાદા છે.
$\Delta = \emptyset$ હોય, તો પહેલી !!{proof}
\Log{G1i}માં પણ !!a{proof} હોય,
પરંતુ બીજી ન હોય, કારણ કે એક સિક્વન્ટની જમણી બાજુ
બંને $!B$ અને~$\lnot !B$ છે.
\sourcecorrection{OLMD-326}{મૂળ અહીં ફરી G1iનો તફાવત માત્ર જમણી બાજુની સંખ્યાની મર્યાદા તરીકે કહે છે. અગાઉના અર્થઘટન-વિભાગમાં નોંધ્યું છે તેમ કેટલાક નિયમોના આધારનું સ્વરૂપ પણ બદલાય છે. અહીં આ મર્યાદા બંને નિષેધ-વૃક્ષો વચ્ચેનો બતાવેલો તફાવત સમજાવે છે.}

\Log{G1c}માં પણ $!A \ident \lforall[x][!B(x)]$
હોય ત્યારે નિયમોનો બીજો ક્રમ વાપરી ન શક્યા હોત તે નોંધો.
આ !!a{proof} \emph{નથી}:
\[
  \AxiomC{}
  \RightLabel{$\pi_1$}
  \Deduce$!B(c) \fCenter !B(c)$
  \RightLabel{\RightR{\lforall}}
  \UnaryInf$!B(c) \fCenter \lforall[x][!B(x)]$
  \RightLabel{\LeftR{\lforall}}
  \UnaryInf$\lforall[x][!B(x)] \fCenter \lforall[x][!B(x)]$
  \DisplayProof
\]
કારણ કે~$\RightR{\lforall}$ની આઇગનચલ-શરત તૂટે છે.

\begin{prop}\ollabel{prop:G3c-axioms}
$!A, \Gamma \Sequent \Delta, !A$ સ્વરૂપનો દરેક સિક્વન્ટ
($!A$ આણ્વિક હોવું \emph{જરૂરી નથી})
\Log{G3c}માં સાબિત થઈ શકે છે.
\end{prop}

\begin{proof}
સાબિતી \olref{prop:G1c-axioms}ની સાબિતી જેવી છે,
પરંતુ આગમન-ધારણા વિશે સાવચેતી લેવી જોઈએ.
$n = \depth{!A} = 0$નો કિસ્સો ફરી સરળ છે.
$n=0$ હોય, તો $!A$ આણ્વિક છે,
અને $!A, \Gamma \Sequent \Delta, !A$ એ~\Log{G3c}નું સ્વયંસિદ્ધ છે.

હવે $n>0$ ધારો. ફરી કિસ્સા જુદા પાડીએ.
આગમન-ધારણા આ છે:
$\depth{!D} < n$ હોય તે બધાં $!D$ માટે,
અને બધા $\Pi$ તથા $\Lambda$ માટે,
$\Log{G3c} \Proves !D, \Pi \Sequent \Lambda, !D$ છે.
\sourcecorrection{OLMD-327}{મૂળ આ સ્થાનધારકને $D!$ લખે છે; બાકીના સૂત્રો પ્રમાણે તેને $!D$ કર્યો છે.}
$!A \ident (!B \land !C)$નો કિસ્સો લઈએ.
આગમન-ધારણા બે વાર વાપરીએ:
એક વાર $!D = !B$, $\Pi = !C, \Gamma$
અને $\Lambda = \Delta$ લઈને
$!B, !C, \Gamma \Sequent \Delta, !B$ની !!a{proof}~$\pi_1$ મેળવીએ;
બીજી વાર $!D = !C$, $\Pi = !B, \Gamma$
અને $\Lambda = \Delta$ લઈને
$!B, !C, \Gamma \Sequent \Delta, !C$ની !!a{proof}~$\pi_2$ મેળવીએ.
\sourcecorrection{OLMD-328}{મૂળ બીજી આગમન-અરજીમાં ડાબા સંદર્ભ માટે જમણા સંદર્ભનું ચિહ્ન Lambda લખે છે. દર્શાવેલા આધાર અને વૃક્ષ મુજબ અહીં એ જ ડાબો Gamma સંદર્ભ રાખ્યો છે.}
$!B \land !C, \Gamma \Sequent \Delta, !B \land !C$ની
!!{proof} આ છે:
\[
\AxiomC{}
\RightLabel{$\pi_1$}
\Deduce$!B, !C, \Gamma \fCenter \Delta, !B$
\RightLabel{\LeftR{\land}}
\UnaryInf$!B \land !C, \Gamma \fCenter \Delta, !B$
\AxiomC{}
\RightLabel{$\pi_2$}
\Deduce$!B, !C, \Gamma \fCenter \Delta, !C$
\RightLabel{\LeftR{\land}}
\UnaryInf$!B \land !C, \Gamma \fCenter \Delta, !C$
\RightLabel{\RightR{\land}}
\BinaryInf$!B \land !C, \Gamma \fCenter \Delta, !B \land !C$
\DisplayProof
\]

$!A \ident \lforall[x][!B(x)]$ માટે
$\lforall[x][!B(x)]$, $\Gamma$ કે~$\Delta$માં
ન આવતો !!a{constant} પસંદ કરો.
$!D = !B(c)$, $\Pi = \lforall[x][!B(x)], \Gamma$
અને $\Lambda = \Delta$ માટે આગમન-ધારણા લગાવો.
તેથી $!B(c), \lforall[x][!B(x)], \Gamma \Sequent \Delta, !B(c)$ની
!!a{proof}~$\pi_1$ મળે છે, જેમાંથી આ મેળવીએ:
\[
\AxiomC{}
\RightLabel{$\pi_1$}
\Deduce$!B(c), \lforall[x][!B(x)], \Gamma \fCenter \Delta, !B(c)$
\RightLabel{\LeftR{\lforall}}
\UnaryInf$\lforall[x][!B(x)], \Gamma \fCenter \Delta, !B(c)$
\RightLabel{\RightR{\lforall}}
\UnaryInf$\lforall[x][!B(x)], \Gamma \fCenter \Delta, \lforall[x][!B(x)]$
\DisplayProof
\]
$c$ની પસંદગીથી $\RightR{\lforall}$ની આઇગનચલ-શરત સંતોષાય છે.
\end{proof}

\begin{prob}
  $!A \ident (!B \lif !C)$ અને $!A \ident \lexists[x][B(x)]$
  હોય તે કિસ્સા કરીને \olref{prop:G3c-axioms}ની સાબિતી આગળ વધારો.
\end{prob}

\end{document}
